\documentclass[12pt,a4paper,oneside]{article}
% nagyon sok kép esetén meggyorsítható a fordítás a draft móddal
% \documentclass[12pt,a4paper,oneside,draft]{report}
% ekkor a képek nem renderelődnek ki, csak placeholder lesz mérethelyesen
\usepackage[utf8]{inputenc} % mindenképp maradjon az utf-8 kódolás
\usepackage[magyar]{babel}
\usepackage[T1]{fontenc}
\usepackage{amsmath}
\usepackage{amsfonts}
\usepackage{amssymb}
\usepackage{graphics} % grafikus elemek, képek berakásához
\usepackage{epsfig} % eps importáláshoz
\usepackage{listings}
\usepackage{sectsty}
\usepackage{enumerate}
\usepackage{lastpage}
\usepackage{setspace}
\usepackage{hyperref} % PDF hivatkozásokhoz kell
\usepackage[hang]{caption}
\usepackage{subcaption}
\usepackage{titling} % a title, author parancsok szabad használatához

\usepackage{acronym}
\usepackage{pdfpages}

% a TikZ rajzoló modul, és a kapcsolási rajz készítő modul, ha kell
\usepackage{tikz}
\usetikzlibrary{shapes.geometric}
%\usepackage{circuitikz}

% az A4 oldal margóinak és méreteinek beállítása
\usepackage[left=25mm,right=25mm,top=20mm,bottom=25mm]{geometry}\pagestyle{plain}

% A sorköz távolság beállítása
% egyszeres sorköz
\singlespacing
% 1,5 sorköz
% \onehalfspacing

% A hivatkozások, és linkek átállítása alapértelmezett színre, fekete-fehér nyomtatáshoz optimalizálva
\hypersetup
{
  	colorlinks,
  	citecolor=black,
 	linkcolor=black,
  	urlcolor=black
}

% --------- A hallgató által kitöltendő rész! --------- %

% a dolgozat típusa
\title{Nagyházi} 

% a dolgozat szerzője
\author{Surányi Lili - , Kozma Dávid Márk - JRY74K, Bordás Máté- }

% a dolgozat típusát adjuk meg tárgyesetben, kisbetűvel (pl.: diplomatervet, szakdolgozatot, témalabor beszámolót stb.)
\newcommand{\dokumentumtipus}{Nagyházi} 

% a dolgozat témája
\newcommand{\dokumentumcim}{Hálózati hibák hatása a QoS-re IP hálózatokban}

% --------- A hallgató által kitöltendő rész! --------- %

\begin{document}
\begin{titlepage}

\begin{figure}
\centering
\includegraphics[width=100mm,keepaspectratio]{logo/bme.pdf}
\end{figure}

\centering
\textbf{Budapesti Műszaki és Gazdaságtudományi Egyetem}\\
\textbf{Villamosmérnöki és Informatikai Kar}\\
\textbf{Hálózati Rendszerek és Szolgáltatások Tanszék}\\
\vspace{5mm}
\includegraphics[width=40mm,keepaspectratio]{logo/HIT_szurke_karikaba_nagy.png}  \\
\vspace{30mm}
\Huge
\dokumentumcim\\
\vspace{54mm}
\Large
\thetitle \\
\vspace{6mm}
\Large
\textbf{\theauthor}\\
\vspace{6mm}
\the\year


\end{titlepage}
\tableofcontents
\include{page/meres}
% végső fordításnál kapcsold ki a draft módot !!

\end{document}